% !TeX root = ../papers/nc-p3.tex
\subsection{三次元接点を持つ有限四辺の完全分類}
\label{sec:maximal-tangent-classification}
三生成部分に点対応がグラフであるという仮定を置かず、残る非被約曲線も判定する。

\begin{lemma}\label{lem:ternary-point-graph-growth}
三生成三関係の標準三次元quadratic Koszul左右AS正則代数 $B=T(H)/(S)$ の
点対応は、平面全体または平面三次曲線上の自己同型のグラフである。三次曲線は非被約でもよい。
\end{lemma}
\begin{proof}
三次交叉を張る $w$ の両端縮約は $H^*\simeq S$ である。
$p\in H^*\setminus0$ での関係の左右評価の階数は等しい。
実際、これらは二次双対の $1$ 次部分から $2$ 次部分への $p$ の左右乗法の階数である。
Frobenius対合で左乗法に随伴なのはNakayama自己同型で移した $p$ の右乗法であり、
この右乗法は元の $p$ の右乗法と一定の基底変更で対応する。
従って左右評価の行列式も同一の $p$ の多項式として定数倍である。

もしこの階数が1なら、右隣の評価像を $ku$、左隣の評価像を $kv$ とし、
三つの位置で順に $H/ku,H/kv,k_p$ を循環する生成空間の商を取る。
$k_p$ へは $p$ で写す。その両側に接する関係は全て零になる。
残る隣接対への三関係の像の次元は高々2である。
実際、$w(p,-,-)$ は $S$ の非零元で、その先頭支持は
残りの末端を縮約して得る $(p\otimes1)S=ku$ だからこの商で消える。
従って三文字の一組で残る商の次元は $2\cdot2-2=2$ 以上となる。
組をつなぐ境界に関係がないので長さ $3n$ で $2^n$ 以上となり、
$H_B(t)=(1-t)^{-3}$ に反する。階数0も両端支持から不可能である。
左右とも各点で階数2以上なので、非零行列式の場合、対応はその三次式の
零点スキーム上で単元二次小行列から復元される唯一の核直線のグラフである。
逆射も右評価から得られる。行列式が恒等零の場合は両射影が平面全体へ同型となる。
\end{proof}

\begin{theorem}\label{thm:maximal-tangent-complete}
標数零の代数閉体上の標準14641型のうち、四辺が有限で、ある点の接空間が三次元のものは
定理\ref{thm:reduced-ternary-graphs}の四種類・各二型で尽きる。
三次元接点はどの対象でも一意である。埋込み次元3という条件は窓の添字保持同型不変量で、
この枝とその補集合の間に同型はない。
\end{theorem}
\begin{proof}
三生成基底は定理\ref{thm:maximal-tangent-ore}で復元され、前補題と
定理\ref{thm:nodal-tangent-windows}により非被約三次曲線の場合だけが残る。
その形は三重直線または二重直線と単直線である。

まず重複直線の台上で $\sigma$ が恒等でないとする。
単直線があればその上でも点ごとの恒等でない。
補題\ref{lem:reduced-ternary-fourth-shift}の固定成分の証明は、その成分の
一般点で曲線が被約でありさえすれば働くからである。
この場合、$\sigma^4=\nu$ は非被約構造上でも成立する。
各成分の一般点の局所環はArtin環であり、$p,\sigma p$ の独立性と
$\det[p,\sigma p,d_i(p,\sigma p)]$ の非零性はその剰余体で成立する。
従って前節までの二平面の証明に使う小行列式は局所環の単元であり、
二つの像は同じ階数2の直和因子となる。両零化直線の一致が
$\sigma^4=\nu$ をこの環上で与える。平面の超曲面には埋込み随伴点がないので、
全一般点の局所環上の等式は全スキーム上の等式である。

二直線なら交点、三重直線なら台の自己同型の固定点を $p$ とする。
$p$ は曲線の重複度3の点で接空間は平面全体である。
定理\ref{thm:reduced-ternary-graphs}の $RD=J$ による微分計算は、
局所方程式の二次項が零であることだけを使うので、この場合も
$d\sigma_p=1$ となる。従って射影変換 $\nu$ は $p$ で恒等微分を持ち、
行列を $I+N$, $N^2=0$ と書ける。$I+tN$ は全 $t$ で曲線を保ち、
制限は $\sigma$ と可換である。実際、整数 $t\ge0$ で既に成立し、
曲線の保存と二つの射の一致は閉条件である。
$\rho=I+N/4$ の曲線への制限を取れば $\tau=\sigma\rho^{-1}$ は位数が4を割り、
$p$ を固定し接空間で恒等である。局所環の各次数商で恒等なので、
各有限次数の近傍で単冪かつ有限位数、従って恒等である。
局所環の分離性から $p$ の局所環で恒等であり、全成分が $p$ を通るので
埋込み随伴点の不在から全曲線上で恒等となる。
これで $\sigma$ は平面に延びるが、非被約な三次超曲面にも消える二次式はなく、
このグラフ上で消える三つの双一次式は平面グラフを定めて矛盾する。

残りは重複直線 $z=0$ の台上で $\sigma=1$ の場合である。
文字を $x,y,z$ とすると、$z=0$ の直線の積への三関係の制限の像は
$k(xy-yx)$ である。像が零なら直線の積全体を含みグラフにならない。
核の二関係の $xz,yz$ 係数は独立である。従属ならある左点 $[x:y:0]$ で
二関係の評価がともに零となり階数2に反する。右側も同様である。
従って $z$ は正規元で、$hz=z\theta(h)$、$\theta(z)=z$ と書ける。
$B/(z)=k[x,y]$ とHilbert級数から $z$ の両側乗法は単射である。
従って $\theta$ は生成元上の可逆一次写像から定まる代数自己同型である。
重複直線の一般点で $z$ は余接の非零な横断方向である。
関係 $h(p)z(\sigma p)=z(p)(\theta h)(\sigma p)$ を $z(p)^2$ を法として読むと、
$z(\sigma p)/z(p)$ は剰余体で非零で、$h=x,y$ に共通である。
ゆえに $\theta$ の $H/kz$ 上の作用は射影直線で恒等、すなわちスカラー $\gamma$ である。
全関係は定数を選んで
\[
 xz=z(\gamma x+a z),\qquad yz=z(\gamma y+b z),\qquad
 xy-yx=z(lx+my+rz)
\]
となる。

$\gamma\ne1$ なら $x,y$ に $z$ の倍を足して $a=b=0$ とする。
第三式の $\theta$ 不変性から $l=m=0$、かつ $(\gamma^2-1)r=0$ を得る。
$r=0$ なら平面グラフなので、倍率変更によって $(\gamma,r)=(-1,1)$ だけが残る。
$\gamma=1$, $a=b=0$ なら $z$ は中心である。$(l,m)\ne(0,0)$ なら
$z\ne0$ 上で対応は対角で、$lx+my+rz=0$ は点ごとに固定される単直線である。
実際、二つの中心関係で次点を消去するとこの一次式がそのまま残る。
これは除外済みであり、残る三重直線型は $l=m=0,r=1$ に取れる
（$r=0$ は平面グラフ）。

最後に $\gamma=1$, $(a,b)\ne(0,0)$ なら $(a,b)=(1,0)$ に直せる。
第三式の $\theta$ 不変性は $l=0$ を与える。
$m=1$ なら $r=0$ は射影移動 $[x:y:z]\mapsto[x-z:y:z]$ の平面グラフであり、
他では $r=1$ にできる。$m\ne1$ なら $y$ への $z$ の倍の加法で $r=0$ とする。
$z\ne0$ での関係からこのとき $y=0$ に単直線があり、移動は $x\mapsto x-z$ である。
$R=yz-zy$, $M_y=xy-yx-mzy$, $M_z=xz-zx-z^2$ に対して三次交叉は
\[
 (x-(m+1)z)R-yM_z+zM_y=R(x+z)+M_yz-M_zy .
\]
従って $\nu[x:y:z]=[x-(m+2)z:y:z]$ となる。この単直線の一般点では
二平面から四移動の等式が得られるので、$m+2=4$、すなわち $m=2$ に限る。

残った四つの基底では行列の係数に自由パラメータはない。
全位置に延びる初期補正の線型空間 $E$ を定理\ref{thm:ternary-ore-linear-space}で計算する。
$P=[u:v:0]$ とし、$g=d(z)(P,P)$、$h=u d(y)(P,P)-v d(x)(P,P)$ と置く。
全ての補正を含む行列を使うと結果は次の通りである。
\[
\begin{array}{c|c|c|c|c|c}
\gamma&a&m&r&\dim E&E\text{ の全域で零}\\ \hline
-1&0&0&1&9&h\\
1&0&0&1&12&g\\
1&1&1&1&9&g\\
1&1&2&0&12&g
\end{array}
\]
再現計算は \path{computations/verify_thick_ternary_exclusion.py} で、
三十行の混合条件から両方向の射影の交叉を反復し、$E$ の全基底を上記の二元斉次式に代入する。
$g$ は二次、$h$ は三次なので一方が恒等零なら必ず射影直線上に共通根がある。
そこで $d(P,P)$ は $P$ に比例し、第四座標は一次元の解を持つ。
これで非被約三次曲線を全て排除した。前節までの各有限表示で指定点が一意であることと、
局所環の接空間が同型で保たれることから残りの主張を得る。
\end{proof}
